\section{Matched direct spline forecasting on OpenDD}
To compare a more complete label roster, we reuse the exposed
OpenDD pedestrian evaluation~\cite{opendd}: \OpenSplineRows{} windows
from \OpenSplineRecordings{} recordings at \OpenSplineSites{} sites,
including \OpenSplineComplete{} fully observed futures. We preserve
eight observations and twelve requested steps at the original stride
of twelve source samples. This is a different corpus and sampling
interval from the moving-drone experiment, not a controlled
missingness intervention.

Each spline fits on the other six sites for each of three seeds,
matching the frozen fixed-head ego EqMotion adaptation~\cite{eqmotion}
on fitting sites, eligible draws, optimizer, robust loss and 800
updates. The new study runs \OpenSplineFits{} fits
(\OpenSplineUpdates{} updates). All three spline conditions use the
same architecture without peers. We compare OLS4 and CV anchoring in
the same centered representation, then change CV history to the
focal comoving representation. The latter also changes the scale
computed by the same past-only preprocessing. No held model selection
or new calibration is performed.

\begin{center}
\small
\resizebox{\columnwidth}{!}{%
\begin{tabular}{lrrr}
\toprule
Forecast & ADE & FDE & Lower gain vs. CV \\
\midrule
CV & \OpenSplineCVADE & \OpenSplineCVFDE & --- \\
OLS4 & \OpenSplineOLSADE & \OpenSplineOLSFDE & \OpenSplineOLSGain \\
Ego EqMotion & \OpenSplineEqADE & \OpenSplineEqFDE & \OpenSplineEqGain \\
OLS4 spline & \OpenSplineOLSSplineADE & \OpenSplineOLSSplineFDE & \OpenSplineOLSSplineGain \\
CV spline & \OpenSplineCVSplineADE & \OpenSplineCVSplineFDE & \OpenSplineCVSplineGain \\
Comoving spline & \OpenSplineMovingADE & \OpenSplineMovingFDE & \OpenSplineMovingGain \\
\bottomrule
\end{tabular}}
\end{center}

ADE pools complete tracks, FDE pools available endpoints, and errors
use publisher coordinates. Gains are CV-normalized pp with 95\%
paired location intervals after averaging seeds. The comoving versus
centered CV spline contrast is \OpenSplineMovingCentered{} pp, and
the comoving spline versus ego EqMotion contrast is \OpenSplineMovingEq{}
pp. Under the original recording checks the comoving spline passes
\OpenSplineRiskPasses{} of \OpenSplineRiskViews{} views, with worst
signed easy degradation of \OpenSplineWorstEasy\%.
This is development evidence, not an independent deployment result.
Omitting the smallest location (\OpenSplineSmallComplete{} complete
windows) in a sensitivity analysis gives \OpenSplineLeaveSmallest{}
pp versus CV; the primary result retains all seven locations.

Keeping \OpenSplineSlots{} intervention recipients, times and peak
displacements identical changes the number of dual-reference passes
from \OpenSplineOldDual{} to \OpenSplineSplineDual{} when the original
shape is replaced by the comoving spline. The paired lower contrast is
\OpenSplineSlotEffect{} pp. An outcome-aware choice between original and
spline yields \OpenSplineSplineCapacity{} feasible views, compared to
\OpenSplineEqCapacity{} for original plus ego EqMotion. Neither bank
resolves the \OpenSplineUndefined{} views without selected easy
reference. The two-choice capacity is not a causal result.
